\section{Design}

\subsection{Architectural sketch}

Below is the architectural sketch of the components structured in the diagram that
highlights the software and hardware components.

\begin{figure}[h]
    \centering
    \includegraphics[width=1\textwidth]{./img/architecturediagram.png}
    \caption{Architectural sketch of the project}
    \label{fig:archi}
\end{figure}

This is the description of each of the elements of the diagram:
\begin{enumerate}
  \item \textbf{ESP32 DevKit microcontroller} -- central processing unit running FreeRTOS
    and executing all three application tasks preemptively
  \item \textbf{DS18B20 temperature sensor (GPIO\,4)} -- digital sensor communicating over
    OneWire; read by TaskAcquisition every 100\,ms in non-blocking mode
  \item \textbf{Green LED (GPIO\,25)} -- output indicator; ON when temperature alert is OFF
  \item \textbf{Red LED (GPIO\,26)} -- output indicator; ON when temperature alert is active
  \item \textbf{LCD 16$\times$2 I2C (SDA=21, SCL=22)} -- display updated by TaskReporter
    every 500\,ms with current temperature and alert state
  \item \textbf{SensorData module} -- shared data store protected by \texttt{xSensorMutex};
    holds temperature value, validity flag, alert state, and debounce counter
  \item \textbf{TaskAcquisition} -- FreeRTOS task (priority 3) that reads the sensor
    periodically and gives \texttt{xSemNewData}
  \item \textbf{TaskConditioning} -- FreeRTOS task (priority 2) blocked on
    \texttt{xSemNewData}; applies hysteresis and debounce, updates LEDs
  \item \textbf{TaskReporter} -- FreeRTOS task (priority 1) that updates LCD and prints
    STDIO reports at 500\,ms intervals
  \item \textbf{GpioDriver (MCAL)} -- wraps all \texttt{digitalRead}/\texttt{digitalWrite}
    calls; only layer that touches Arduino hardware APIs directly
\end{enumerate}

\subsection{Electrical sketch}

The circuit diagram outlines the setup for this project, showing the connections between
the ESP32, two LEDs (green, red), the DS18B20 temperature sensor, the LCD I2C module,
and the computer for serial communication. Each LED is connected to its respective GPIO
pin through a $220\,\Omega$ current-limiting resistor. The DS18B20 data line (GPIO\,4)
requires a $4.7\,\mathrm{k\Omega}$ pull-up resistor to 3.3\,V for reliable OneWire
communication.

\begin{figure}[h]
    \centering
    \includegraphics[width=0.4\textwidth]{./img/circuit_sketch_wokwi.jpg}
    \caption{Circuit sketched on wokwi}
    \label{fig:wokwi}
\end{figure}

Figure \ref{fig:wokwi} details the connections and components that make up the system.
This visual documentation is crucial for replicating the setup or diagnosing issues, as
it clearly shows how each component is integrated into the overall design.

\subsection{Flow chart}

The Figure \ref{fig:flow} illustrates the FreeRTOS task interaction flow. TaskAcquisition
runs every 100\,ms, reads the DS18B20, updates shared data under the mutex, and gives the
binary semaphore. TaskConditioning unblocks immediately, reads the shared temperature,
applies the \texttt{ProcessChannel} hysteresis + debounce logic, and updates the alert
state and LEDs. TaskReporter runs every 500\,ms independently and reads shared data to
refresh the LCD and serial output.

\begin{figure}[h]
    \centering
    \includegraphics[width=0.8\textwidth]{./img/flow.png}
    \caption{Flow chart of the algorithm}
    \label{fig:flow}
\end{figure}

\clearpage
\subsection{Code design}

Code is developed following the principle of separating the interface from the
implementation, by creating for each header file \texttt{.h} a respective \texttt{.cpp}
file. The project uses \textbf{PlatformIO} with three library directories:

\begin{enumerate}
  \item \textbf{MCAL layer} (\texttt{lib/MCAL/GpioDriver}) -- hardware-specific code that
    wraps \texttt{Arduino.h} GPIO calls. All direct hardware access is encapsulated here.
  \item \textbf{ECAL layer} (\texttt{lib/ECAL/}) -- electronic component drivers:
    \texttt{TempSensor} (DS18B20 via OneWire + DallasTemperature),
    \texttt{LedDriver} (two LEDs),
    \texttt{LcdDisplay} (LiquidCrystal\_I2C), and
    \texttt{SerialIO} (UART initialization).
  \item \textbf{SRV layer} (\texttt{lib/SRV/}) -- the \texttt{SensorData} shared-state
    module and the three FreeRTOS task implementations.
\end{enumerate}

The \texttt{SensorData} module centralizes all shared data behind mutex-protected
getters and setters, and owns the FreeRTOS synchronization primitives:

\begin{lstlisting}
// Declared in SensorData.h, initialized in SensorData_Init()
extern SemaphoreHandle_t xSemNewData;   // binary semaphore: new data available
extern SemaphoreHandle_t xSensorMutex; // mutex: protects all shared sensor fields
\end{lstlisting}

The signal conditioning logic is encapsulated in a single reusable function:
\begin{lstlisting}
static AlertState_t ProcessChannel(
    float value, float threshHigh, float threshLow,
    AlertState_t currentAlert, uint8_t *debounce)
{
    // Hysteresis
    AlertState_t desired;
    if      (value >= threshHigh) desired = ALERT_ON;
    else if (value <= threshLow)  desired = ALERT_OFF;
    else                          desired = currentAlert; // no change in dead band

    // Debounce: require DEBOUNCE_COUNT consecutive confirmations
    if (desired != currentAlert) {
        (*debounce)++;
        if (*debounce >= DEBOUNCE_COUNT) {
            *debounce = 0;
            return desired;  // transition confirmed
        }
    } else {
        *debounce = 0;       // reset counter if condition does not persist
    }
    return currentAlert;     // no state change yet
}
\end{lstlisting}

TaskReporter uses \texttt{vTaskDelayUntil} for drift-corrected 500\,ms periodicity
and alternates between a compact CSV line (every cycle) and a full structured report
(every 5\,seconds, i.e. every 10th cycle):
\begin{lstlisting}
// CSV plotter line (every 500 ms)
printf("%.1f,%.1f,%.1f\n", temp, TEMP_THRESH_HIGH, TEMP_THRESH_LOW);

// Full structured report (every 5 s)
if ((reportCounter++ % FULL_REPORT_EVERY) == 0) {
    printf("\n--- SENSOR REPORT (t=%lu ms) ---\n",
           xTaskGetTickCount() * portTICK_PERIOD_MS);
    printf("Temp  : %.1f C (valid=%d, alert=%s, dbg=%d)\n",
           temp, tempOk, (tAlert ? "ON" : "OFF"), tDbg);
    printf("Thresh: ON >= %.1f C, OFF <= %.1f C\n",
           TEMP_THRESH_HIGH, TEMP_THRESH_LOW);
    printf("--------------------------------\n\n");
}
\end{lstlisting}
